\section{Introduction}\label{sec:intro}

Quantum computations are intrinsically susceptible to noise and hence, quantum error correction (QEC) is essential for achieving fault tolerant quantum computing. The current landscape of the field is changing as it is transitioning from the noisy intermediate scale era towards more concrete demonstrations of fault tolerance. The current paradigm relies on qubit-based codes and architectures but an alternative approach has been proposed in the avenue of qudits. This paper will specifically look at the $q=3$ case of qudits, referred to as qutrits, and will work in the early fault tolerant quantum computing (early-FTQC) regime. At this stage, the goal is to realize small logical qutrits with protection against errors, before full scale QEC is feasible. 

Qutrits are three-level systems that present compelling alternatives for early FTQC research, where the aim is to study encoded logical devices before full error-correction cycles are practical. They display a larger and more connected computational state space than qubits: a single qutrit stores $\log_2 3 \approx 1.585$ bits of quantum information, increasing information density and, in principle, reducing both computational and ancillary overheads. Algorithmic primitives admit direct ternary advantages, for example, base-3 decompositions reduce \emph{(i)} circuit depth and \emph{(ii)} the count of costly entangling/non-Clifford gates in Shor’s modular exponentiation and generalized-Toffoli subroutines \cite{Gokhale2019}. Moreover, implementing the quantum Fourier transform (QFT) directly in \emph{base-3} reduces the register length from \(q_2=\lceil \log_2 N\rceil\) bits to \(q_3=\lceil \log_3 N\rceil\) trits; since a QFT over \(q\) digits uses \(\Theta(q^2)\) controlled rotations / controlled-\(\mathrm{SUM}\) operations, the entangling-gate count (and leading-order depth) drops by a factor \((\log_3 N/\log_2 N)^2 \approx 0.40\) at fixed \(N\), with explicit qudit-QFT circuit constructions given in \cite{PavlidisFloratos2021}. Similarly, Grover's search algorithm implemented with qutrits uses multi-level oracles and diffusion operators, reducing the gate count and circuit depth per iteration on hardware with native qutrit control \cite{Ivanov2012}. Beyond algorithms, higher-dimensional encodings improve quantum communication: in qutrit quantum key distribution (QKD), each carrier can encode more than one bit of key and protocols exhibit higher error tolerance than binary schemes \cite{BechmannPasquinucciPeres2000,Cerf2002}. Higher-dimensional systems strengthen foundational certificates used in quantum communication. With three-level systems, entangled pairs achieve larger violations of the Collins–Gisin–Linden–Massar–Popescu (CGLMP) Bell inequality than any qubit-based test, which increases the noise margin available to device-independent protocols \cite{Collins2002}. Further, qutrits are the smallest systems that can violate the Klyachko–Can–Binicioglu–Shumovsky (KCBS) contextuality inequality, improving robustness in semi-device-independent protocols \cite{Klyachko2008}. Qutrits also provide a natural language for simulating spin-1 and \(\mathrm{SU}(3)\) models and related phenomena, broadening the space of near-term scientific targets \cite{Senko2015}.
From a QEC standpoint, \emph{magic-state distillation} (MSD) in odd prime dimensions is especially attractive: quantum Reed Muller codes provide transversal qutrit gates enabling small block, high rate MSD protocols across all prime \(q\) \cite{CampbellPRX2012}, and working directly with \((q>2)\)-level systems further improves MSD overheads and provides routes to universal non-Clifford gates \cite{CampbellPRL2014}. Despite these advantages, systematic stabilizer-code design and logical benchmarking for qudits remain comparatively underexplored and merit systematic study. These considerations align with hardware reality: superconducting transmons, trapped ions, and photonic platforms natively provide multilevel structure.  Multi-platform demonstrations of coherent qutrit control, randomized benchmarking, and entangling operations now exist \cite{Morvan2021,Ringbauer2022,Goss2024}.

In this work, we employ the $[[5,1,2]]_3$ quantum error detection code to encode logical information. The code layout is tailored to the connectivity of IQM’s superconducting Garnet QPU, ensuring native compatibility with the hardware topology. We rigorously characterize the code properties and its constant-depth transversal logical Clifford gates using the \texttt{sdim} qudit stabilizer simulator \cite{sdim}, establishing a numerical baseline prior to hardware execution. Our primary objective is to benchmark the average fidelity of logical Clifford gates and demonstrate a quantifiable advantage over their physical counterparts. To achieve this, we adapt standard protocols to implement Logical Randomized Benchmarking with Detection (LRB-D). Since the distance $d=2$ code supports error detection but not faithful correction, LRB-D isolates logical performance by post-selecting on error-free trajectories. Crucially, our analysis highlights that logical dynamics can exhibit emergent non-Markovian behavior even under Markovian physical noise, a phenomenon driven by the coupling to syndrome subspaces \cite{Ziyad2025}. We demonstrate that LRB-D, particularly through the use of intermediate post-selection, effectively suppresses these non-Markovian memory effects by filtering leakage trajectories, thereby recovering a single exponential decay and validating the logical performance gain.